% =====================================================================
\chapter{Sumatorios, productorios y binomio de Newton}
\label{tema:4}
% =====================================================================

La probabilidad de una secuencia es un producto de probabilidades condicionadas, su logaritmo es una suma, y la pérdida de entrenamiento es una media sobre millones de tokens. Manejar con soltura la notación $\sum$ y $\prod$ es, por tanto, imprescindible. Este tema sigue el capítulo~2 «Sums» de Graham, Knuth y Patashnik~\cite{gkp} y el capítulo~14 «Sums and Asymptotics» de Lehman, Leighton y Meyer~\cite{mcs}; el binomio de Newton sigue el §9.6 de OpenStax~\cite{openstax}.

\section{Notación y principio de inducción}

\begin{definicion}[Sumatorio y productorio]
Dados enteros $m\le n$ y números $a_m,\dots,a_n$,
\[
\sum_{k=m}^{n}a_k=a_m+a_{m+1}+\dots+a_n,\qquad \prod_{k=m}^{n}a_k=a_m\,a_{m+1}\cdots a_n.
\]
La letra $k$ es un \emph{índice mudo}: puede sustituirse por cualquier otra sin cambiar el valor. Si $n<m$, la suma es \emph{vacía} y vale~0, y el producto vacío vale~1. Más en general, $\sum_{k\in S}a_k$ denota la suma sobre un conjunto finito $S$ de índices. El \emph{factorial} es $n!=\prod_{k=1}^{n}k$, con $0!=1$.
\end{definicion}

Formalmente, la suma se define por recurrencia: $\sum_{k=m}^{m-1}a_k=0$ y $\sum_{k=m}^{n}a_k=\sum_{k=m}^{n-1}a_k+a_n$. Por eso las fórmulas para sumas se demuestran de forma natural por inducción.

\begin{teorema}[Principio de inducción]\label{teo:induccion}
Sea $P(n)$ una afirmación sobre los enteros $n\ge n_0$. Si $P(n_0)$ es cierta y, para cada $n\ge n_0$, $P(n)$ implica $P(n+1)$, entonces $P(n)$ es cierta para todo $n\ge n_0$.
\end{teorema}

Este principio es equivalente al principio del buen orden (todo subconjunto no vacío de $\N$ tiene mínimo); ambos se estudian en~\cite[cap.~2 y~5]{mcs} y~\cite[cap.~10]{hammack}.

\section{Propiedades de los sumatorios}

\begin{proposicion}\label{prop:sumas}
Para números $a_k,b_k$ y una constante $c$:
\begin{enumerate}
  \item linealidad: $\sum_{k=m}^n(ca_k+b_k)=c\sum_{k=m}^na_k+\sum_{k=m}^nb_k$;
  \item partición: si $m\le p<n$, $\sum_{k=m}^na_k=\sum_{k=m}^pa_k+\sum_{k=p+1}^na_k$;
  \item suma de una constante: $\sum_{k=m}^nc=(n-m+1)\,c$;
  \item cambio de índice: $\sum_{k=m}^na_k=\sum_{j=m+r}^{n+r}a_{j-r}$ para todo entero $r$, y $\sum_{k=0}^na_k=\sum_{k=0}^na_{n-k}$.
\end{enumerate}
\end{proposicion}

No hay reglas análogas para el producto: en general $\sum a_kb_k\neq\bigl(\sum a_k\bigr)\bigl(\sum b_k\bigr)$.

\section{Sumas clásicas}

\begin{teorema}[Suma aritmética]\label{teo:aritmetica}
Para todo $n\in\N$, $\displaystyle\sum_{k=1}^nk=\frac{n(n+1)}2$. En general, si $a_k=a_1+(k-1)d$, entonces $\sum_{k=1}^na_k=\dfrac{n(a_1+a_n)}{2}$.
\end{teorema}

\begin{teorema}[Suma geométrica]\label{teo:geometrica}
Para $r\neq1$ y $n\in\N_0$, $\displaystyle\sum_{k=0}^nr^k=\frac{1-r^{n+1}}{1-r}$. Si $r=1$, la suma vale $n+1$.
\end{teorema}

\begin{proposicion}[Sumas telescópicas]\label{prop:telescopica}
Para cualesquiera números $a_m,\dots,a_{n+1}$,
\[
\sum_{k=m}^n(a_{k+1}-a_k)=a_{n+1}-a_m.
\]
\end{proposicion}

\addtocounter{definicion}{1}%

\begin{teorema}[Sumas de potencias]\label{teo:potencias}
Para todo $n\in\N$,
\[
\sum_{k=1}^nk^2=\frac{n(n+1)(2n+1)}6,\qquad\sum_{k=1}^nk^3=\Bigl(\frac{n(n+1)}2\Bigr)^2.
\]
\end{teorema}

\addtocounter{definicion}{1}%

\section{Sumas dobles}

Una suma doble $\sum_{i}\sum_{j}a_{ij}$ suma los elementos de una tabla. Si los índices recorren un rectángulo, el orden de suma puede intercambiarse:
\[
\sum_{i=1}^m\sum_{j=1}^na_{ij}=\sum_{j=1}^n\sum_{i=1}^ma_{ij}.
\]
Si recorren un triángulo, hay que reescribir los límites~\cite[§2.4]{gkp}:
\[
\sum_{i=1}^n\sum_{j=i}^na_{ij}=\sum_{1\le i\le j\le n}a_{ij}=\sum_{j=1}^n\sum_{i=1}^ja_{ij}.
\]

\begin{proposicion}\label{prop:producto-sumas}
$\displaystyle\Bigl(\sum_{i=1}^ma_i\Bigr)\Bigl(\sum_{j=1}^nb_j\Bigr)=\sum_{i=1}^m\sum_{j=1}^na_ib_j$. En particular, $\displaystyle\Bigl(\sum_{i=1}^na_i\Bigr)^2=\sum_{i=1}^na_i^2+2\sum_{1\le i<j\le n}a_ia_j$.
\end{proposicion}

\addtocounter{definicion}{1}%

\section{Productorios y binomio de Newton}

Los productorios cumplen $\prod_k(a_kb_k)=\prod_ka_k\prod_kb_k$, $\prod_{k=1}^nc=c^n$ y, si todos los $a_k$ son positivos, $\ln\prod_ka_k=\sum_k\ln a_k$ (\cref{tema:5}).

\begin{definicion}[Números combinatorios]
Para enteros $0\le k\le n$, $\displaystyle\binom nk=\frac{n!}{k!\,(n-k)!}$, y $\binom nk=0$ si $k<0$ o $k>n$.
\end{definicion}

\begin{lema}[Identidad de Pascal]\label{lem:pascal}
Para $1\le k\le n$, $\displaystyle\binom nk+\binom n{k-1}=\binom{n+1}k$.
\end{lema}

\begin{teorema}[Binomio de Newton]\label{teo:binomio}
Para todos $a,b\in\R$ y $n\in\N_0$, $\displaystyle(a+b)^n=\sum_{k=0}^n\binom nka^{n-k}b^k$.
\end{teorema}

\begin{corolario}\label{cor:binomio}
$\displaystyle\sum_{k=0}^n\binom nk=2^n$ y, si $n\ge1$, $\displaystyle\sum_{k=0}^n(-1)^k\binom nk=0$.
\end{corolario}

\addtocounter{definicion}{1}%

\addtocounter{definicion}{1}%

